\documentclass[12pt, letterpaper]{article}
\usepackage[margin=1in]{geometry}
\usepackage{setspace}
\usepackage{amsmath}
\usepackage{graphicx}
\usepackage{times}
\usepackage{float}
\usepackage{booktabs}

\doublespacing

\begin{document}

\begin{center}
    \textbf{Siamese Transformer Networks for Edit Distance Embedding and Evaluation} \\
    Hsin Ya Lin \\
    Institute of Communications, National Yang Ming Chiao Tung University
\end{center}

\section*{Abstract}
This report evaluates the efficacy of a Siamese Transformer architecture in mapping error correction code operations into a geometric embedding space. Initial experiments demonstrated that generative decoding struggles with edit distance tasks due to shift-induced error propagation. By utilizing an encoder-only architecture trained with Triplet Margin Loss, the network successfully learns the structural relationships between sequences. The model achieves an RMSE of 2.092 and a Triplet Accuracy of 0.7238, demonstrating that attention mechanisms can effectively map discrete sequence mutations into a measurable metric space, outperforming standard recurrent and convolutional baselines in absolute error and separation accuracy.

\section*{Methodology and Architecture}
Standard generative sequence-to-sequence models face fundamental limitations when applied to edit distance problems. A single uncorrected insertion or deletion operation shifts the predicted sequence, causing positional loss functions to heavily penalize structurally valid outputs. To resolve this instability, the architecture was refactored into an encoder-only Siamese network designed to evaluate Euclidean distances in a continuous manifold.

The dataset utilizes a triplet schema consisting of anchor, positive (correlated), and negative (corrupted) sequences. The architecture is a Siamese Transformer featuring an embedding dimension of $d_{model} = 256$, utilizing 8 attention heads across 4 encoder layers. Mean pooling collapses the high-dimensional encoder memory into a single representative embedding vector. 

The network was trained over 30 epochs using the AdamW optimizer. Optimization was governed by a Triplet Margin Loss function, defined as:
\begin{equation}
    \mathcal{L}(a, p, n) = \max \{ d(a_i, p_i) - d(a_i, n_i) + {\rm margin}, 0 \}
\end{equation}
where $d$ is the pairwise Euclidean distance, $a_i$ is the anchor, $p_i$ is the positive sample, $n_i$ is the negative sample, and the margin is set to 1.0. To calculate scale-comparable evaluation metrics (RMSE and MAE), a 1D linear regression was fitted to project the continuous Euclidean embedding distances onto the discrete Levenshtein edit distance scale.

\section*{Empirical Results}
Performance was evaluated by calculating the distances for 10,000 sequence triplets. The network yielded an average positive distance of 2.7128 and an average negative distance of 3.3750, resulting in a Triplet Accuracy of 0.7238.

\begin{figure}[H]
    \centering
    \includegraphics[width=0.8\textwidth]{plots/distance_histogram.png}
    \caption{Distribution of Euclidean distances in the 256-dimensional embedding space, demonstrating structural separation between correlated and corrupted sequence pairs.}
    \label{fig:histogram}
\end{figure}

Figure \ref{fig:histogram} visualizes the mathematical separation between the mean distributions, confirming the network pulled correlated sequences together while pushing corrupted sequences apart.

\begin{figure}[H]
    \centering
    \includegraphics[width=0.8\textwidth]{plots/distance_scatter.png}
    \caption{Correlation mapping between true Levenshtein edit distance and the predicted Euclidean embedding distance.}
    \label{fig:scatter}
\end{figure}

\begin{figure}[H]
    \centering
    \includegraphics[width=0.8\textwidth]{plots/pca_clusters.png}
    \caption{Principal Component Analysis projecting the 256-dimensional Siamese embedding manifold into a 2D space, colored by sequence corruption severity.}
    \label{fig:pca}
\end{figure}

The Principal Component Analysis in Figure \ref{fig:pca} illustrates that the attention layers capture the underlying variance of the edit distance manifold, organizing sequences spatially based on structural integrity.

\section*{Discussion and Baseline Comparison}
To contextualize the performance, the Siamese Transformer was benchmarked against established architectural baseline data. Table \ref{tab:baselines} outlines the comparative metrics across the evaluated models.

\begin{table}[H]
\centering
\caption{Comparative Edit Distance Metrics Across Network Architectures}
\vspace{0.2cm}
\begin{tabular}{lcccccc}
\toprule
\textbf{Model} & \textbf{RMSE} $\downarrow$ & \textbf{MAE} $\downarrow$ & \textbf{Pearson} $\uparrow$ & \textbf{Spearman} $\uparrow$ & \textbf{Triplet Acc.} $\uparrow$ & \textbf{Hamming Loss} $\downarrow$ \\ 
\midrule
Transformer (Gen.) & 3.772 & 2.919 & 0.311 & 0.223 & 0.530 & 0.373 \\ 
GRU & 3.244 & 2.530 & 0.592 & 0.575 & 0.636 & 0.310 \\ 
CNN-ED & 2.906 & 2.263 & 0.621 & 0.612 & 0.684 & 0.325 \\ 
CGK & 2.622 & 2.062 & 0.688 & 0.693 & 0.721 & N/A \\ 
\textbf{Siamese (Proposed)} & \textbf{2.092} & \textbf{1.670} & \textbf{0.455} & \textbf{0.452} & \textbf{0.7238} & \textbf{N/A} \\ 
Specialized Lab Arch. & 1.541 & 1.180 & 0.912 & 0.914 & 0.835 & 0.049 \\ 
\bottomrule
\end{tabular}
\label{tab:baselines}
\end{table}

The data reveals a significant architectural tradeoff. The proposed Siamese Transformer achieves an RMSE of 2.092 and an MAE of 1.670, substantially outperforming traditional models like CGK and CNN-ED in minimizing absolute distance error. Furthermore, its Triplet Accuracy of 0.7238 proves a strong capacity for binary sequence separation. 

However, the correlation coefficients (Pearson 0.455, Spearman 0.452) trail the convolutional and alignment baselines. This indicates that while Triplet Margin Loss effectively establishes tight absolute boundaries and pushes local pairs past the defined margin, it does not inherently force a globally linear ranking across the entire severity scale. Because the architecture maps relationships rather than predicting discrete sequences, Hamming Loss is structurally inapplicable. While the specialized lab architecture maintains the highest overall fidelity, the variance captured by the Siamese encoder proves that multi-head attention mechanisms provide a highly accurate estimation of error correction metrics without the overhead of generative decoding.

\end{document}